\section{Threats to Validity}
\label{sec:threats}

\paragraph{Internal validity.} A difference between two configurations can come from the oracle, not the database:
ties under \texttt{ORDER BY}, unordered \texttt{LIMIT}, \texttt{any\_value}, row identifiers, floating-point noise near
zero, or a time-zone dependence. Each new generator's first hits were usually such oracle errors. We reduce this risk by
the generator contract (unique sort keys), result normalization, replay in a fresh process, minimization that keeps the
verdict, and a person who decides which side is right by reading the rule's code before filing. Upstream confirmation
is the final check: 11 of 16 issues are labelled reproduced and 4 are fixed. The method that found each bug is taken
from our notes and run logs; for the issues found before the automated loop existed, the attribution is ours and was
not measured. One early bug that we first believed new had been fixed upstream the same day (PR \#28405); we caught it
before filing because triage checks the newest development build, but such races can recur.

\paragraph{External validity.} We tested one database, DuckDB, and mostly one release, 2.0.0, which was new during
our project. A new major release probably contains more optimizer bugs than a mature one, so our yield may not transfer
to older or more heavily tested systems. Most bugs belong to one family, so the results say little about other bug
classes such as concurrency or storage bugs. The LLMs we used change over time, and a later model may behave
differently on the same prompts.

\paragraph{Construct validity.} We count \emph{filed issues}, not confirmed root causes. An issue that maintainers
later close as expected behaviour would reduce the count; none has been so far. Coverage numbers come from one round of
Method C and measure lines, which predict bug finding only weakly. Reach is measured by a plan change, which shows that
the rewrite fired but not that it fired in the interesting way. Our cost measure counts LLM calls, not tokens or wall
time, because the runs were interrupted by usage limits and by the machine sleeping.
